% Part: proof-theory
% Chapter: natural-deduction
% Section: translation-N2i

\documentclass[../../../include/open-logic-section]{subfiles}

\begin{document}

\olfileid{pt}{nat}{tni}

\olsection{له \Log{N2} څخه \Log{G2} ته ژباړه}

\begin{prop}\ollabel{prop:N2-to-G2}
که $\Log{N2c}\ (\Log{N2i}) \Proves \Gamma \Sequent !A$ وي، نو
$\Log{G2c}\ (\Log{G2i}) + \Cut \Proves \Gamma' \Sequent \Delta$،
چې $\Gamma'$ د !!{formula}s د تکرار شمېر لرونکے هغه سټ دے
چې له $\Gamma$ څخه د نښو په لرې کولو جوړېږي؛ که $!A$ د
$\lfalse$ فارمول نۀ وي $\Delta = \{!A\}$، او که وي
$\Delta = \emptyset$۔
\textbf{سرچينه‌يي سمون OLCMP-159:} د حدسي ژباړې پاييز نظام
G2i دے؛ د سرچينې نقل شوے N2i نوم سم شو۔
\end{prop}

\begin{proof}
ثبوت بيا د $\Gamma \Sequent !A$ د !!{proof}~$\delta$ پر
لوړوالي استقرا ده۔ که $\delta$ د $x:!A \Sequent !A$
بديهيت وي، $\pi$ اړوند بديهيت $!A \Sequent !A$ دے، يا
$\lfalse \Sequent \ $ که $!A \ident \lfalse$ وي۔

که $\delta$ په استنتاجي قاعده ختمېږي، حالتونه بېلوو:
\begin{enumerate}
  \item که $R$ د \Intro{\lif} قاعده وي، اصلي فارمول ئې
  $!B \lif !C$ وي، او $x:!B$ د مقدمې په سياق کښې موجود خو
  د پايلې په سياق کښې نۀ وي، نو $\delta$ داسې ختمېږي:
  \[
  \AxiomC{}
  \RightLabel{$\delta_1$}
  \Deduce$x:!B, \Gamma \fCenter !C$
  \RightLabel{\Intro{\lif}}
  \UnaryInf$\Gamma \fCenter !B \lif !C$
  \DisplayProof
  \]
  د استقرا فرضيه د $!B, \Gamma' \Sequent !C$
  !!a{proof}~$\pi_1$ راکوي، يا د $!B, \Gamma' \Sequent \ $
  که $!C \ident \lfalse$ وي۔ له دې څخه $\pi$ د
  \RightR{\lif} په کارولو اخلو؛ که $!C \ident \lfalse$ وي،
  لومړے \RightR{\Weakening} کاروو۔
  \item که $R$ د \Intro{\lif} قاعده وي، خو $x:!B$ د
  مقدمې په سياق کښې نۀ وي، نو $\delta$ داسې ختمېږي:
  \[
  \AxiomC{}
  \RightLabel{$\delta_1$}
  \Deduce$\Gamma \fCenter !C$
  \RightLabel{\Intro{\lif}}
  \UnaryInf$\Gamma \fCenter !B \lif !C$
  \DisplayProof
  \]
  د استقرا فرضيه د $\Gamma' \Sequent !C$
  !!a{proof}~$\pi_1$ هغه وخت راکوي چې پاييز فارمول دروغ
  نۀ وي۔ که دروغ وي، د استقرا ثبوت تش ښے اړخ لري؛ لومړے
  ښے اړخ کمزوري کوو چې هماغه دروغ فارمول ښي اړخ ته راشي۔
  بيا په دواړو حالتونو کښې $\pi$ د~$!B$ سره د
  \LeftR{\Weakening} او ورپسې د~\RightR{\lif} په کارولو اخلو۔
  \textbf{سرچينه‌يي سمون OLCMP-160:} په بې‌فرضه ايستلو
  حالت کښې هم د دروغ پاييز فارمول تش ښے اړخ د شرطي داخلولو
  مخکې بېرته په فارمول اړول پکار دي؛ سرچينې دا حالت پرېښودے و۔
  \item که $R$ د \Elim{\lif} قاعده وي، $\delta$ داسې ختمېږي:
  \[
  \AxiomC{}
  \RightLabel{$\delta_1$}
  \Deduce$\Gamma_1 \fCenter !B \lif !A$
  \AxiomC{}
  \RightLabel{$\delta_2$}
  \Deduce$\Gamma_2 \fCenter !B$
  \RightLabel{\Elim{\lif}}
  \BinaryInf$\Gamma_1, \Gamma_2 \fCenter !A$
  \DisplayProof
  \]
  چې $\Gamma = \Gamma_1 \cup \Gamma_2$ دے۔ د استقرا فرضيه
  د $\Gamma_1' \Sequent !B \lif !A$ !!{proof}s $\pi_1$ او
  د $\Gamma_2' \Sequent !B$ ثبوت $\pi_2$ راکوي، په دې
  وروستي بيان کښې هغه وخت چې کوچنۍ مقدمه دروغ نۀ وي۔ هغه
  حالت چې کوچنۍ مقدمه دروغ وي، لاندې بېل بيانېږي۔ کله چې
  پاييز فارمول هم دروغ نۀ وي، !!{proof}~$\pi$ داسې اخلو:
  \[
  \AxiomC{}
  \RightLabel{$\pi_1$}
  \Deduce$\Gamma_1' \fCenter !B \lif !A$
  \AxiomC{}
  \RightLabel{$\pi_2$}
  \Deduce$\Gamma_2' \fCenter !B$
  \Axiom$!A \fCenter !A$
  \RightLabel{\LeftR{\lif}}
  \BinaryInf$!B \lif !A, \Gamma_2' \fCenter !A$
  \RightLabel{\Cut}
  \BinaryInf$\Gamma_1', \Gamma_2' \fCenter !A$
  \DisplayProof
  \]
  که پاييز فارمول دروغ وي خو کوچنۍ مقدمه دروغ نۀ وي، په
  همدې رسم کښې د هويت بديهيت پر ځاے د دروغ-تش سېکوېنټ
  بديهيت اخلو؛ د شرطي کيڼې قاعدې او ورپسې د پرېکون پايلې
  هم تش ښے اړخ لري۔
  \textbf{سرچينه‌يي سمون OLCMP-162:} د قضیې د دروغ پاييز
  فارمول ژباړه تش ښے اړخ غواړي؛ عمومي رسم په همدغه حالت
  کښې د دروغ-تش بديهيت سره لوستل کېږي۔
  د تکرار شمېر لرونکے سټ $\Gamma_1' \cup \Gamma_2'$ ښايي
  له $\Gamma' = (\Gamma_1 \cup \Gamma_2)'$ سره برابر نۀ وي۔
  مثلاً، که دواړه $\Gamma_1$ او $\Gamma_2$ د $x:!C$ فارمول
  ولري، نو $\Gamma_1' \cup \Gamma_2'$ د $!C$ دوه نقلونه لري،
  خو $(\Gamma_1 \cup \Gamma_2)'$ يوازې يو نقل لري۔ د
  \LeftR{\Contraction} د مناسبو استنتاجونو په زياتولو دا سمولے شو۔

  که $!B \ident \lfalse$ وي، پر $\delta_2$ د استقرا
  فرضيه د $\Gamma_2' \Sequent \ $ !!a{proof} راکوي۔ له دې
  څخه !!{proof}~$\pi$ د \LeftR{\Weakening} د مناسبو
  استنتاجونو په زياتولو جوړوو چې د لوې مقدمې اړين سياق
  هم ولري؛ بيا د~$!A$ سره \RightR{\Weakening} يوازې هغه
  وخت کاروو چې پاييز فارمول دروغ نۀ وي۔ که پاييز فارمول
  دروغ وي، ښے اړخ هماغه تش پرېږدو۔
  \textbf{سرچينه‌يي سمون OLCMP-161:} دويم فرعي ثبوت د
  کوچنۍ مقدمې فارمول ثابتوي؛ د تش ښي اړخ شرط ئې د کوچنۍ
  مقدمې دروغ‌والے دے، نۀ د وروستۍ پايلې دروغ‌والے۔ د
  وروستۍ دروغ پايلې دپاره ښے اړخ تش ساتل کېږي۔
\item که $R$ د \FalseCl قاعده وي، $\delta$ د ښودل شوي
  نښه لرونکي فرض د موجوديت په حالت کښې داسې ختمېږي:
  \[
  \AxiomC{}
  \RightLabel{$\delta_1$}
  \Deduce$x: \lnot !A, \Gamma \fCenter \lfalse$
  \RightLabel{\FalseCl}
  \UnaryInf$\Gamma \fCenter !A$
  \DisplayProof
  \]
  د استقرا له فرضيې د $\lnot !A, \Gamma' \Sequent \ $
  !!a{proof} $\pi_1$ شته۔ که د ايستلو نښه لرونکے فرض په
  اصل کښې موجود نۀ وي، لومړے کيڼ اړخ کمزوري کوو چې نفي
  شوې مقدمه ورزياته شي۔
  \textbf{سرچينه‌يي سمون OLCMP-164:} کلاسيکي قاعده د
  فرض له ايستلو پرته هم کارېږي؛ سرچينې يوازې د موجود فرض رسم ښودلے و،
  د بې‌فرضه ايستلو حالت هم په دغه کيڼ کمزوري کولو رانغاړو۔
  !!{proof}~$\pi$ داسې اخلو:
  \[
  \Axiom$!A \fCenter !A$
  \RightLabel{\RightR{\lnot}}
  \UnaryInf$\fCenter !A, \lnot !A$
  \AxiomC{}
  \RightLabel{$\pi_1$}
  \Deduce$\lnot !A, \Gamma' \fCenter $
  \RightLabel{\Cut}
  \BinaryInf$\Gamma' \fCenter !A$
  \DisplayProof
  \]
  \textbf{سرچينه‌يي سمون OLCMP-163:} د G2 په سېکوېنټ کښې
  د فرض نښه نۀ پاتې کېږي؛ د استقرا په مقدمه کښې پاتې اضافي ليبل لرې شو۔
  که پاييز فارمول دروغ وي، د دې رسم په پاے کښې له
  دروغ-تش بديهيت سره يو بل پرېکون کوو، څو ښے اړخ د
  قضیې له شرط سره سم تش شي۔
  \textbf{سرچينه‌يي سمون OLCMP-165:} د کلاسيکي وروستي
  حالت دروغ پايله هم بايد تش ښي اړخ ته واړول شي؛ سرچينې
  د خپلې قضیې دغه شرط دلته نۀ و پوره کړے۔
\end{enumerate}
پام وکړئ چې يوازې د \FalseCl{} ژباړه داسې
!!a{proof}~$\pi$ ورکوي چې ښي اړخ ته له يوه !!{formula}
څخه زيات ولري۔ يعنې د~\Log{N2i} د !!{proof} ژباړه يو
\Log{G2i}-!!{proof} دے۔
\textbf{د ثبوت د بشپړتيا سرچينه‌يي يادونه OLCMP-166:}
سرچينه يوازې د شرطي داخلولو دوه حالتونه، د شرطي ايستلو
او د کلاسيکي دروغ قاعدې حالتونه ورکوي؛ نور منطقي او
کميت‌لرونکي قاعدوي حالتونه نۀ دي تشريح شوي۔ د هغو د
ورک استدلال له امله دا ورکړل شوې برخه د بشپړ ثبوت په
توګه نۀ تصديقېږي۔ قضیه ناسمه نۀ بلل کېږي او نوے ورک
ثبوت يا حل شوے تمرين دلته نۀ دے زيات شوے۔
\end{proof}

\end{document}
